% (c) Danny Lee Eldridge · AEGIS Community License · CC BY-NC-SA 4.0
\documentclass[aps,prd,twocolumn,superscriptaddress,nofootinbib,showkeys]{revtex4-2}
\usepackage{amsmath,amssymb,hyperref}
\begin{document}
\title{The Reverse Thread: The Inflationary Species Plateau Reused as Thawing Dark Energy, Scored Against DESI DR2, with a Measured Tuning Window}
\author{Danny Eldridge}
\email{dannyeldridge54@gmail.com}
\affiliation{Independent researcher, Oklahoma City, Oklahoma, USA}
\date{19 November 2026}
\begin{abstract}
In a companion paper~\cite{P1}, inflation is driven by a species-scale plateau $V\propto1-Ae^{-\beta\phi}$ with $A=16\zeta(2)/3\zeta(3)=7.30$ and $\beta=\sqrt{D-2}$. Here we postulate that a second, slow field carries the same plateau at the dark-energy scale. The field is frozen through most of cosmic history and begins to roll only recently: a ``reverse'' of inflation. This is an instance of quintessential $\alpha$-attractors~\cite{L15,AKLV,BMS}, with $\alpha=2/(3\beta^2)=1/12,1/6,1/3$ for $D=10,6,4$. We integrate the background with matter and radiation and project $w(a)$ onto $(w_0,w_a)$. We then score the result against DESI DR2 combined with the CMB and each of Pantheon+, Union3 and DESY5~\cite{DESI}. The score uses a Gaussian approximation whose correlation is calibrated so that $\Lambda$CDM reproduces DESI's own $\Delta\chi^2$ values, and Section~II lists its limits.
(i) All three tilts reach $\Delta\chi^2\le6.18$ for all three supernova sets simultaneously, against $10.7$, $17.4$ and $21.0$ for $\Lambda$CDM. The common best point is $w_0\simeq-0.79$, $w_a\simeq-0.36$.
(ii) This happens only in a narrow window of starting heights, $0.8$--$1.5\%$ of the viable range in $\ln A e^{-\beta\phi_i}$, just before the field would roll into $V<0$ before today.
(iii) The height $V_\Lambda\sim10^{-122}M_P^4$ is put in by hand.
The model therefore fits, but with two coincidences. Its sharp predictions are $w\ge-1$ at all times and, if DESI's central values hold, $(w_0,w_a)\approx(-0.79,-0.36)$, to be scored on DESI DR3 and Euclid. All computations were pre-registered.
\end{abstract}
\keywords{dark energy; quintessence; $\alpha$-attractors; species scale; DESI; inflation; fine-tuning; pre-registration}
\maketitle

\section{Model}
The species function $E(\tau,\tfrac32)=2\zeta(3)y^{3/2}+4\zeta(2)y^{-1/2}+\ldots$ gives the inflationary plateau $1-Ae^{-\beta_D\psi}$ of Ref.~\cite{P1}, with $A=7.30$. As a postulate (P-G), we give a second slow field $\phi$ the same shape at a tiny height:
\begin{equation}
V(\phi)=V_\Lambda\left(1-A\,e^{-\beta\phi}\right),\qquad \beta\in\{2\sqrt2,\,2,\,\sqrt2\}.
\end{equation}
We work in units $M_P=H_0=1$, with $\Omega_m=0.31$ and $\Omega_r=9\times10^{-5}$. The field starts frozen at $a=10^{-3}$, and $V_\Lambda$ is fixed by shooting to $H(a{=}1)=1$. The single free parameter is the starting height, $\mathrm{amp}\equiv Ae^{-\beta\phi_i}$. Starting points that cannot reach $H_0$ without rolling into $V<0$ are recorded as non-viable. We fit $w=w_0+w_a(1-a)$ by least squares over $a\in[1/3,1]$. When the field reaches its minimum it oscillates and redshifts like matter, but this lies in the future here and has no consequence today.

The leading plateau term of an E-model $\alpha$-attractor, $(1-e^{-\sqrt{2/3\alpha}\,\phi})^2$, has the same form, so $\beta=\sqrt{2/(3\alpha)}$. Borys, Masias and Scalisi~\cite{BMS} already showed that quintessential $\alpha$-attractors fit DESI DR2 for $\alpha=\mathcal O(1)$. Our values $\alpha=1/12$--$1/3$ lie inside their preferred ranges. They also noted a ``double fine-tuning'' of today's potential value and slope. What is specific to this note is the coefficient $A$, the tilts taken from the inflation sector, and a quantitative measure of the tuning.

\section{Scoring and its limits}
We use DESI DR2 + CMB + SN, $w_0w_a$CDM~\cite{DESI}: Pantheon+ $(-0.838\pm0.055,\,-0.62^{+0.22}_{-0.19})$, Union3 $(-0.667\pm0.088,\,-1.09^{+0.31}_{-0.27})$, DESY5 $(-0.752\pm0.057,\,-0.86^{+0.23}_{-0.20})$. We take a two-dimensional Gaussian with asymmetric $w_a$ errors. Its correlation $\rho$ is chosen so that the $\Lambda$CDM point $(-1,0)$ has DESI's quoted $\Delta\chi^2_{\rm MAP}=10.7$, $17.4$ and $21.0$, giving $\rho=-0.55$, $-0.54$ and $-0.59$.
This is \emph{not} a fit to the BAO, CMB and SN likelihoods. $\Omega_m$ is not refit, and our CPL projection differs from DESI's procedure. The scores are indicative only, and a full-likelihood analysis is the needed next step. As a cross-check, our improvement over $\Lambda$CDM (13.6 for Union3, 15.1 for DESY5) is of the same size as the thawing-class gain reported in Ref.~\cite{QP} ($\approx13.5$, against 17.2 for CPL).

\section{Results}
\begin{table}[b]\caption{Fine grid, 300 starting heights per tilt. ``Window'' is the fraction of the viable $\ln(\mathrm{amp})$ range in which $\Delta\chi^2\le6.18$ for all three SN sets at once. The last column is the smallest worst-case $\Delta\chi^2$ over the three sets, with its $(w_0,w_a)$.}\label{tab:t}
\begin{ruledtabular}\begin{tabular}{lccc}
$\beta$ ($\alpha$) & window & best worst-case & $(w_0,w_a)$\\
$2\sqrt2$ (1/12) & 1.5\% & 5.21 & $(-0.78,-0.38)$\\
$2$ (1/6) & 0.9\% & 5.59 & $(-0.79,-0.36)$\\
$\sqrt2$ (1/3) & 0.8\% & 5.69 & $(-0.79,-0.35)$\\
\end{tabular}\end{ruledtabular}\end{table}
Table~\ref{tab:t} summarises the results. Every tilt passes all three supernova sets in a narrow window, and all three converge on $(w_0,w_a)\approx(-0.79,-0.36)$. The window sits immediately before the non-viable region, so the field is caught just as it leaves the plateau. Away from the window the model drifts back towards $\Lambda$CDM's poor score (the field still frozen) or becomes non-viable. An independent integration in the scale factor $a$ with a different solver reproduces the best point to all printed digits.

\section{What is assumed, and what would falsify it}
We assume the following: postulate P-G (a slow field carries the inflationary plateau shape); the height $V_\Lambda$, put in by hand; a starting height tuned to about $\pm2\%$; and the Gaussian approximation of Sec.~II. The model was formed after DESI DR2 and must be scored on future data only. It is excluded if DR3 or Euclid require $w<-1$ at any redshift, or if they keep $w_a\lesssim-0.6$ at $w_0\gtrsim-0.8$, which this thawing plateau cannot reach.

\emph{Integrity note.} Locks: PREREG\_GENESIS\_DE, \_DE2 and \_DE3 (D1626--D1628, repository branch \texttt{unify-bg}). Fixes for two crashes and one output error were logged before reruns. Two corrections are on record. D1626a: a maximum $|w_a|$ had been read from a subsampled printout. D1627: an assumed correlation $\rho=-0.9$ was replaced by the calibration above, and this moved the verdict for Union3 from tension to consistency.

\begin{thebibliography}{99}
\bibitem{P1} D.~Eldridge, ``Inflation from the running of the species scale'', doi:10.5281/zenodo.24378141; flavour companion doi:10.5281/zenodo.24379357.
\bibitem{L15} E.~V.~Linder, ``Dark energy from $\alpha$-attractors'', arXiv:1505.00815.
\bibitem{AKLV} Y.~Akrami, R.~Kallosh, A.~Linde, V.~Vardanyan, ``Dark energy, $\alpha$-attractors, and large-scale structure surveys'', arXiv:1712.09693.
\bibitem{BMS} A.~Borys, J.~Masias, M.~Scalisi, ``Quintessential $\alpha$-attractors fit DESI'', arXiv:2606.02701.
\bibitem{DESI} DESI Collaboration, DESI DR2 Results II, arXiv:2503.14738, Tables V--VI.
\bibitem{QP} ``Quintessence and phantoms in light of DESI 2025'', arXiv:2506.21542.
\end{thebibliography}
\end{document}
